% Vista editorial de corpus/source/masas/02_ley_general.tex, líneas 363--491: cuerpo exacto salvo el retítulo académico del epígrafe sobre el ciclo de 108 estados. Las líneas 1--263 ya residen exactamente en 72_rev2_nucleo_ley_i.tex y las líneas 264--362 en 72_rev2_algoritmo_universal.tex.
\chapter{Ampliación sistemática de la teoría de masas}
\Needspace{7\baselineskip}
\section{Fundamento unificado de la masa}
\Needspace{7\baselineskip}
\subsection{El objeto anterior a la magnitud}
La teoría comienza antes de la cifra. Una masa físicamente individualizada no es un número real acompañado por el nombre de una partícula, sino el término final de una biografía matemática. APP aporta el espacio finito de incidencias y las dos operaciones correlativas; el TRIT conserva la orientación temporal de cada transición; el TPK transporta fase, hoja, acarreo, frontera, ruta, memoria y supervivencia. La Mecánica Dimensional convierte después esa biografía en una sección de la recta de masa. El orden completo es

\[
\mathrm{APP}\longrightarrow\mathrm{TRIT}\longrightarrow\mathrm{TPK}
\longrightarrow\Gamma^{\rm surv}
\longrightarrow\gamma_{\rm obs}
\longrightarrow\mathscr L_\gamma
\longrightarrow\sigma(\gamma)
\longrightarrow\chi_{\rm mass}
\longrightarrow\mathcal T_{90,120}
\longrightarrow\mathcal L_M.
\]
Cada flecha retiene información que la siguiente utiliza. La celda no reemplaza la ruta; la ruta no reemplaza el registro; una firma no reemplaza su registro genealógico; un carácter no reemplaza la sección dimensional sobre la que actúa. Esta separación explica por qué dos partículas pueden compartir carga o espín y, sin embargo, poseer masas diferentes: su biografía enriquecida no es la misma. También explica por qué partícula y antipartícula comparten masa: la involución cambia la orientación impar y conserva la memoria par que desciende al cociente de masa.

APP consta de nueve por nueve posiciones publicadas. Cada posición admite cuatro orientaciones iniciales, de modo que el dominio orientado contiene

\[
81\cdot4=324
\]
rutas candidatas. La clasificación interna produce 56 familias de ruta y 13 multisecciones. Estas cardinalidades son simultáneas: 81 cuenta celdas, 324 cuenta orientaciones, 56 cuenta cocientes por historia y 13 cuenta tipos operatorios gruesos. Ninguna tabla de partículas sustituye a esos dominios.

\Needspace{7\baselineskip}
\subsection{Unidad discreta del continuo que sostiene la masa}
La recta de masa no se adosa a cinco teorías independientes del continuo. Procede del único estado APP--TRIT--TPK del que emergen correlativamente los aspectos espectral, solenoidal, de calibre, coherente y determinantal, designados en la notación de construcción por \(sp,sol,gau,coh,det\). Esos símbolos son discriminantes de lectura, no subdominios previos. Si \(\mathcal C_{\rm disc}\) denota la estructura común, la flecha generadora tiene la forma

\[
\mathfrak E_{\rm cont}:\mathcal C_{\rm disc}
\longrightarrow
\bigl(\mathcal C_\infty;
\mathcal O_{sp},\mathcal O_{sol},\mathcal O_{gau},
\mathcal O_{coh},\mathcal O_{det}\bigr),
\]
donde las cinco observaciones pertenecen a un mismo objeto límite, comparten fibras, orientación, acarreo, memoria, frontera y lectores, y se transportan por los mismos cuadrados de refinamiento. La masa utiliza su lectura determinantal para la escala de acción, su lectura espectral para el carácter y las torres, su lectura de calibre para los sectores nulos y cargados, y sus lecturas coherente y solenoidal para la propagación y la composición. Ninguna de ellas se introduce suponiendo de antemano que el estado pertenece a una clase física clásica.

Esta unidad es una condición de tipado: separar uno de esos aspectos y conservar los otros produciría un objeto distinto, incapaz de sostener a la vez ruta, acción, carga, amplitud y masa. La presente teoría no vuelve a demostrar el continuo completo; incorpora su construcción unitaria como fundamento y verifica que ninguna de las inserciones de masas lo disgrega.

\Needspace{7\baselineskip}
\subsection{Escala absoluta y cocientes}
El lector determinantal local tiene forma normal de Smith

\[
\operatorname{SNF}(H_4)=\operatorname{diag}(1,2,2,4),
\]
de donde

\[
|\operatorname{coker}H_4|=16,
\qquad
\Sigma_H=\varphi_{\rm HMT}\alpha_{\rm HMT}^{16},
\qquad
\operatorname{ord}_{10}(\Sigma_H)=-34.
\]
La década se obtiene antes de la carta SI. Las secciones pre y retorno de la recta de acción son

\[
\hbar_{\rm pre}=H_5\,10^{-34}u_S,
\qquad
\hbar_{\rm ret}=\eta_{\rm ret}^{(\deg)}\,10^{-34}u_S,
\]
y su cociente bidireccional es

\[
R_{\rm act}=\frac{H_5}{\eta_{\rm ret}^{(\deg)}}.
\]
En la rama de reloj elegimos explícitamente como sección de acción la sección de retorno, \(\hbar_{\rm int}:=\hbar_{\rm ret}\). Un reloj declarado \(t_0\) proporciona

\[
\omega_0=\frac{2\pi}{108t_0},
\qquad
\varepsilon_{\rm clk}=\hbar_{\rm ret}\omega_0,
\qquad
m_{\rm clk}=\varepsilon_{\rm clk}c_{\rm int}^{-2}.
\]
Adoptar \(u_M^{\rm clk}:=m_{\rm clk}\) como origen de la carta es una normalización explícita. No identifica automáticamente esta sección con una base energética o másica previamente declarada. El carácter produce primero el operador basal

\[
\widehat M_X^{(0)}
=\mathfrak e_0m_{\rm clk}\,
\chi_{\rm full}(\sigma_X-\sigma_e,\eta_X-\eta_e),
\]
y el lector bidireccional actúa después. Para cada presentación \(r\in\{D,\Omega\}\), sea \(\widehat B_X^r=(\widehat B_X^r)^*\) sobre la fibra finita; en una realización infinita se exige además un dominio común denso para el cálculo funcional. Entonces

\[
\widehat M_X^{\beta,r}
=R_{\rm act}^{\widehat B_X^r/2}
\widehat M_X^{(0)}
R_{\rm act}^{\widehat B_X^r/2}.
\]
En una fibra escalar, \(\widehat B_X^r=\kappa_XI\), de modo que

\[
m_X^\beta=m_X^{(0)}R_{\rm act}^{\kappa_X}.
\]
Si las coordenadas de torre ya se expresan relativamente al electrón, se omite \(-\eta_e\). Al formar un cociente, la sección dimensional común se cancela, pero no la diferencia de lectura bidireccional:

\[
\frac{m_Y^\beta}{m_X^\beta}
=\chi_{\rm full}(\sigma_Y-\sigma_X,\eta_Y-\eta_X)
R_{\rm act}^{\kappa_Y-\kappa_X}.
\]
Por tanto, \(10^{-34}\) fija la escala absoluta de acción, energía y masa; no es un corrector relativo oculto. La información fina de una especie sólo puede entrar mediante su carácter, su reloj relativo o el cociente orientado de las secciones pre y retorno.

\Needspace{7\baselineskip}
\subsection{Cuanto interno de la recta de masa inducido por el ciclo de 108 estados}
La autodualidad Planck--CT108 proporciona, en la normalización interna

\[
\hbar_{\rm int}=c_{\rm int}=t_0=1,
\]
la sección

\[
m_{108}^{\rm int}
=0.058177641733144319230789692282953757\ldots .
\]
Este número es un cuanto normalizado de la recta interna de masa. No es la masa de una partícula, no lleva por sí solo una carta MeV o GeV y no puede insertarse como fila adicional de especies. Su función es comprobar la compatibilidad entre la carta de acción, el reloj CT108 y la autodualidad de la recta de masa. La transición a una masa nominal exige todavía la ruta, la firma, el operador y la sección dimensional de la especie.

\Needspace{7\baselineskip}
\subsection{Existencia de la salida HMT–MD para extensiones supervivientes y criterio de reducción escalar por rango de fibra}
Fijada una extensión superviviente, un lector celular, una firma entera, una familia de coeficientes de torre absolutamente convergente y una sección dimensional de masa, existe una salida HMT--MD bien tipada. Si la fibra posee un único punto fijo compatible con la involución, la salida se reduce a un escalar canónico. Si la fibra tiene rango mayor, la salida primaria es el operador espectral sobre la fibra; una masa escalar física requiere el funcional de acoplamiento que individualiza la representación observada.

Este teorema no degrada las evaluaciones escalares recuperadas. Distingue tres afirmaciones: el operador interno está construido; la evaluación de una firma declarada es exacta; la determinación prospectiva de esa firma es un teorema adicional cuando no se deduce del propio lector. La comparación experimental se efectúa únicamente después de fijar cuál de esas tres afirmaciones se publica.

\Needspace{7\baselineskip}
